\section{ARQUITECTURA Y DISEÑO DE FIRMWARE}

\subsection{ARQUITECTURA DE SOFTWARE DEL SSEP}

El firmware del SSEP.2.1 (SSW.2.1) está escrito en C y organizado en tres capas. Cada
capa sólo usa a la inmediata inferior, de modo que la lógica de la aplicación trabaja con
magnitudes y estados (RPM, centésimas de grado, mando por mil, modo) y no conoce pines,
registros ni periféricos. El diagrama de ESW por capa se incluye en el Anexo B.

\begin{itemize}
  \item \textbf{Capa de aplicación} (\texttt{Aplicacion/}): el programa principal y la
  lógica de decisión (modos, alarma y regulador PI). No incluye ningún encabezado de
  dispositivos; esta regla se comprueba automáticamente en el build.
  \item \textbf{Capa de servicios} (\texttt{Servicios/}): adapta los dispositivos a
  magnitudes útiles para la aplicación (velocidad, temperatura filtrada, mando
  normalizado, consignas y telemetría) y es el único camino de la aplicación hacia el
  hardware.
  \item \textbf{Capa de dispositivos} (\texttt{Dispositivos/}): un módulo por elemento
  físico. Configura los pines y periféricos y los accede mediante la biblioteca de drivers
  LPC17xx original, que no se modifica. Incluye además un soporte de plataforma
  (\texttt{Dispositivos/soporte/}) con el arranque del reloj y la adquisición ADC/DMA
  compartida por el potenciómetro y el LM35.
\end{itemize}

Los parámetros del sistema se concentran en encabezados de \texttt{Configuracion/}:
límites de consigna y mando (\texttt{Limites.h}), pulsos por vuelta
(\texttt{Medicion.h}), ganancias y período del regulador (\texttt{Control.h}), umbral e
histéresis de la alarma (\texttt{Alarma.h}) y tensión de referencia del ADC
(\texttt{Analogica.h}). Todos pueden redefinirse al compilar y se validan con
\texttt{\#error} si quedan fuera de rango.

\subsubsection{MODELO DE EJECUCIÓN}

Se utiliza un \textbf{lazo cooperativo} en \texttt{main} con tres interrupciones breves.
Las ISR sólo mueven datos y actualizan contadores; todas las decisiones se toman en el lazo
principal. No se usa memoria dinámica ni números en punto flotante.

\begin{longtable}{|p{3.6cm}|p{1.6cm}|p{9.4cm}|}
\hline
\textbf{ISR} & \textbf{Prioridad} & \textbf{Función} \\ \hline
\endhead
\texttt{TIMER0\_IRQHandler} & 1 & Cada 500\,ms lee el contador del sensor óptico, publica
los pulsos de la ventana y reprograma MR0. \\ \hline
\texttt{UART3\_IRQHandler} & 2 & Pasa los bytes recibidos de la FIFO al buffer circular de
recepción y carga la FIFO de transmisión desde el buffer circular de envío. \\ \hline
\texttt{DMA\_IRQHandler} & 3 & Cuenta cada par de muestras ADC transferido y marca que
hay datos disponibles. \\ \hline
\end{longtable}

Los datos compartidos entre ISR y lazo principal tienen un único escritor por variable
(buffers circulares de un productor y un consumidor, contadores de 32 bits), por lo que no
se requieren secciones críticas. La lectura de la ventana de medición relee el contador de
ventanas para detectar si la interrupción intervino durante la consulta.

La llegada de una ventana nueva de medición marca el \textbf{período de control de
500\,ms}: alarma, modo, mando y telemetría se ejecutan una vez por ventana, en ese orden
(ver sección 2.1). El regulador PI usa ese mismo período.

\subsubsection{SECUENCIA DE INICIALIZACIÓN}

\begin{enumerate}
  \item \texttt{Sistema\_Init}: tiempo de acceso a flash, PLL0 a 100\,MHz y arranque de
  ADC/DMA; espera acotada hasta tener el primer par de muestras.
  \item \texttt{Driver\_Motor\_Init} y mando cero.
  \item \texttt{Driver\_Alarmas\_Init}; la alarma queda apagada, o encendida si falló el
  paso 1 (en ese caso el programa se detiene aquí).
  \item \texttt{Mandos\_Init} (llave selectora) y \texttt{Comunicacion\_Init} (UART3).
  \item Modo inicial según la llave.
  \item \texttt{Medicion\_Velocidad\_Init}: contador del sensor y TIMER0. Se hace al
  final para que la primera ventana de 500\,ms empiece con todo listo.
\end{enumerate}

\subsubsection{HERRAMIENTAS Y VERIFICACIÓN}

El proyecto se compila con MCUXpresso IDE (proyecto en la raíz del repositorio) y también
con CMake y GCC Arm (\texttt{arm-none-eabi}), sin errores ni advertencias. Además se
compilan los mismos fuentes para la PC con dispositivos simulados, lo que permite ejecutar
pruebas unitarias y de integración de todos los servicios y de la aplicación, también con
\textit{sanitizers} de memoria. Todas las funciones están documentadas con Doxygen
(\texttt{@file}, \texttt{@brief}, \texttt{@details}, \texttt{@note}) e identificadas con su
ESW y USW, y la generación de la documentación no produce advertencias.

\subsection{FIRMWARE DEL SSEP}

\subsubsection{DESCRIPCIÓN DE USW DE LA CAPA DE APLICACIÓN}

\textbf{ESW.2.1.1 -- main.c / main.h}

\begin{longtable}{|p{2.4cm}|p{4.6cm}|p{7.6cm}|}
\hline
\textbf{USW} & \textbf{Función} & \textbf{Descripción} \\ \hline
\endhead
USW.2.1.1.1 & \texttt{main} & Inicializa y repite el lazo principal. Si la
inicialización falla queda detenido con el motor en cero y la alarma encendida. \\ \hline
USW.2.1.1.2 & \texttt{Aplicacion\_Inicializar} & Ejecuta la secuencia de inicialización
de la sección anterior. Devuelve \texttt{false} si no arrancó la adquisición
analógica. \\ \hline
USW.2.1.1.3 & \texttt{Aplicacion\_Iterar} & Una vuelta del lazo, sin esperas: en MANUAL
aplica el potenciómetro al motor; si hay una medición nueva ejecuta el período de
control. \\ \hline
USW.2.1.1.4 & \texttt{Aplicacion\_CalcularMando} & Mando del período según el estado:
potenciómetro en MANUAL, salida del PI en AUTO\_ACTIVO y cero en AUTO\_ESPERA. \\ \hline
USW.2.1.1.5 & \texttt{Aplicacion\_Periodo} & Período de 500\,ms: alarma, modo (descarta
la recepción pendiente al pasar de MANUAL a automático), mando y telemetría. \\ \hline
\end{longtable}

\textbf{ESW.2.1.2 -- Procesamiento.c / Procesamiento.h}

Contiene toda la lógica de decisión, sin acceso a periféricos ni llamadas a servicios:
\texttt{main} le entrega las lecturas y aplica sus resultados.

\begin{longtable}{|p{2.4cm}|p{4.6cm}|p{7.6cm}|}
\hline
\textbf{USW} & \textbf{Función} & \textbf{Descripción} \\ \hline
\endhead
USW.2.1.2.11 & \texttt{Procesamiento\_\-InicializarModo} & Fija el estado de arranque:
MANUAL o AUTO\_ESPERA según la llave. Borra consigna e integral. \\ \hline
USW.2.1.2.12 & \texttt{Procesamiento\_\-ActualizarModo} & Aplica las transiciones de
estado de la sección 2.2. Una consigna mayor a 185\,RPM se ignora sin saturarla. \\ \hline
USW.2.1.2.13 & \texttt{Procesamiento\_\-Consigna} & Devuelve la consigna vigente (cero fuera
de AUTO\_ACTIVO). \\ \hline
USW.2.1.2.14 & \texttt{Procesamiento\_\-EvaluarAlarma} & Alarma con histéresis: se
enciende por encima de 45,00\,\textdegree C y se apaga a 43,00\,\textdegree C o menos; entre
ambos valores conserva la decisión anterior. \\ \hline
USW.2.1.2.9 & \texttt{Procesamiento\_\-Regular} & Un paso del regulador PI (ver
abajo). \\ \hline
\end{longtable}

\textbf{Regulador PI.} El error se calcula en milésimas de RPM, $e = 1000 \cdot
\mathrm{consigna} - \mathrm{velocidad}$, y el mando es
\[
u = K_p \, e + I, \qquad 0 \leq u \leq 1000
\]
redondeado una sola vez al entero más cercano. Luego la integral se actualiza como $I
\leftarrow I + K_i \, e$, limitada a $[0, 1000]$. Como medida \textit{anti-windup}, la
integral no se actualiza si el mando sin limitar está saturado en el sentido del error.
Las ganancias se expresan como fracciones enteras en \texttt{Control.h} ($K_p =
4/1000$ y $K_i = 1/1000$ por milésima de RPM, con $K_i$ ya discretizada para el paso de
500\,ms) y se calculan en aritmética entera de 64 bits, sin desborde posible. Con consigna
cero el mando es cero y la integral se borra; la integral también se borra en MANUAL, de
modo que cada entrada al modo automático regula desde cero.
\pendiente{actualizar las ganancias si se ajustaron durante la puesta en marcha y agregar
el resultado de los ensayos de respuesta}.

\subsubsection{DESCRIPCIÓN DE USW DE LA CAPA DE SERVICIOS}

\begin{longtable}{|p{2.4cm}|p{4.6cm}|p{7.6cm}|}
\hline
\textbf{USW} & \textbf{Función} & \textbf{Descripción} \\ \hline
\endhead
\multicolumn{3}{|l|}{\textbf{ESW.2.1.3 -- Medicion\_Velocidad}} \\ \hline
USW.2.1.3.3 & \texttt{Medicion\_Velocidad\_\-Init} & Arranca el contador del sensor óptico
y TIMER0; la primera ventana empieza en esta llamada. \\ \hline
USW.2.1.3.4 & \texttt{Medicion\_Velocidad\_\-Actualizar} & Si cerró una ventana nueva,
devuelve la velocidad en milésimas de RPM: $\mathrm{pulsos} \cdot 60\,000 \cdot 1000 / (60
\cdot 500)$, en 64 bits. Cada \texttt{true} marca el período de control. No recorta a
185\,RPM. \\ \hline
\multicolumn{3}{|l|}{\textbf{ESW.2.1.4 -- Temperatura}} \\ \hline
USW.2.1.4.3 & \texttt{Temperatura\_Leer} & Toma una lectura del LM35 y devuelve el
promedio móvil de las últimas cuatro, en centésimas de grado (filtro de 2\,s). \\ \hline
\multicolumn{3}{|l|}{\textbf{ESW.2.1.5 -- Mandos}} \\ \hline
USW.2.1.5.6 & \texttt{Mandos\_Init} & Configura la entrada de la llave selectora. \\ \hline
USW.2.1.5.7 & \texttt{Mandos\_LeerManual} & Convierte la muestra del potenciómetro (0 a
4095) en mando por mil: $(N \cdot 1000 + 2047) / 4095$. \\ \hline
USW.2.1.5.8 & \texttt{Mandos\_LeerModo} & Devuelve MANUAL (llave abierta) o AUTOMÁTICO
(llave cerrada). \\ \hline
\multicolumn{3}{|l|}{\textbf{ESW.2.1.6 -- Driver\_Alarmas}} \\ \hline
USW.2.1.6.2 & \texttt{Driver\_Alarmas\_Init} & Configura LED y buzzer, apagados. \\ \hline
USW.2.1.6.3 & \texttt{Driver\_Alarmas\_\-Aplicar} & Enciende o apaga juntos el LED rojo y el
buzzer; cada dispositivo resuelve su polaridad. \\ \hline
\multicolumn{3}{|l|}{\textbf{ESW.2.1.7 -- Comunicacion}} \\ \hline
USW.2.1.7.11 & \texttt{Comunicacion\_Init} & Configura UART3 y vacía la línea en
recepción. \\ \hline
USW.2.1.7.12 & \texttt{Comunicacion\_\-RecibirConsigna} & Arma la línea recibida hasta el
LF y entrega la próxima consigna válida; descarta líneas inválidas o largas. \\ \hline
USW.2.1.7.13 & \texttt{Comunicacion\_\-Descartar} & Olvida lo recibido; se usa al entrar en
automático. \\ \hline
USW.2.1.7.14 & \texttt{Comunicacion\_\-EnviarTelemetria} & Arma la trama
\texttt{RPM=...;TEMP=...;ESTADO=...} y la encola para transmitir, sin esperar. \\ \hline
USW.2.1.7.15 & \texttt{Comunicacion\_\-Decodificar} & (Interna) Valida el formato
\texttt{RPM n} y convierte el valor. \\ \hline
USW.2.1.7.8 a .10 & \texttt{AgregarTexto}, \texttt{AgregarEntero}, \texttt{AgregarFijo} &
(Internas) Escriben literales, enteros y números con punto decimal fijo en la trama, sin
biblioteca estándar ni punto flotante. \\ \hline
\multicolumn{3}{|l|}{\textbf{ESW.2.1.8 -- Driver\_Motor}} \\ \hline
USW.2.1.8.3 & \texttt{Driver\_Motor\_Init} & Configura el PWM del motor con mando
cero. \\ \hline
USW.2.1.8.4 & \texttt{Driver\_Motor\_Aplicar} & Aplica un mando de 0 a 1000 por mil;
valores mayores se limitan a 1000. \\ \hline
\multicolumn{3}{|l|}{\textbf{Soporte de plataforma -- Sistema}} \\ \hline
-- & \texttt{Sistema\_Init} & Configura el reloj, arranca ADC/DMA y espera, con un límite
de consultas, el primer par de muestras. \\ \hline
\end{longtable}

\subsubsection{DESCRIPCIÓN DE USW DE LA CAPA DE DISPOSITIVOS}

\begin{longtable}{|p{2.4cm}|p{4.6cm}|p{7.6cm}|}
\hline
\textbf{USW} & \textbf{Función} & \textbf{Descripción} \\ \hline
\endhead
\multicolumn{3}{|l|}{\textbf{ESW.2.1.9 -- Motor}} \\ \hline
USW.2.1.9.2 & \texttt{Motor\_Init} & PWM1 a 1\,kHz con MR0 = 1000, canal 2 en P2.1, mando
inicial cero. \\ \hline
USW.2.1.9.1 & \texttt{Motor\_Write} & Escribe el mando directamente en MR2 (un paso por
unidad por mil); rechaza valores mayores a 1000. \\ \hline
\multicolumn{3}{|l|}{\textbf{ESW.2.1.10 -- Sensor\_Optico}} \\ \hline
USW.2.1.10.2 & \texttt{SensorOptico\_Init} & TIMER1 como contador de flancos ascendentes
en CAP1.0 (P1.18), con pull-up. \\ \hline
USW.2.1.10.1 & \texttt{SensorOptico\_Read} & Devuelve el conteo acumulado de
pulsos. \\ \hline
\multicolumn{3}{|l|}{\textbf{ESW.2.1.11 -- LM35}} \\ \hline
USW.2.1.11.3 & \texttt{LM35\_Read} & Lee la muestra transferida por DMA y devuelve la
temperatura en centésimas de grado. \\ \hline
USW.2.1.11.4 & \texttt{LM35\_Convert} & (Interna) Conversión entera con redondeo al más
cercano según VREFP y 10\,mV/\textdegree C. \\ \hline
\multicolumn{3}{|l|}{\textbf{ESW.2.1.12 -- Potenciometro}} \\ \hline
USW.2.1.12.1 & \texttt{Pot\_Read} & Devuelve el resultado de 12 bits (0 a 4095) de la
muestra transferida por DMA. \\ \hline
\multicolumn{3}{|l|}{\textbf{ESW.2.1.13 -- Selector}} \\ \hline
USW.2.1.13.2 & \texttt{Selector\_Init} & P2.6 como entrada con pull-up. \\ \hline
USW.2.1.13.1 & \texttt{Selector\_Read} & Lee el nivel de P2.6. \\ \hline
\multicolumn{3}{|l|}{\textbf{ESW.2.1.14 -- LED}} \\ \hline
USW.2.1.14.2 & \texttt{LED\_Init} & P0.22 como salida, LED apagado. \\ \hline
USW.2.1.14.1 & \texttt{LED\_Write} & Enciende (nivel bajo) o apaga el LED rojo. \\ \hline
\multicolumn{3}{|l|}{\textbf{ESW.2.1.15 -- Buzzer}} \\ \hline
USW.2.1.15.2 & \texttt{Buzzer\_Init} & P2.7 como salida, buzzer apagado. \\ \hline
USW.2.1.15.1 & \texttt{Buzzer\_Write} & Activa (nivel alto) o desactiva el buzzer. \\ \hline
\multicolumn{3}{|l|}{\textbf{ESW.2.1.16 -- UART}} \\ \hline
USW.2.1.16.3 & \texttt{UART3\_Init} & UART3 a 9600 8N1 en P0.0/P0.1, FIFO e
interrupciones de recepción y transmisión. \\ \hline
USW.2.1.16.1 & \texttt{UART\_Read} & Retira bytes del buffer circular de recepción (64
bytes), sin esperar. \\ \hline
USW.2.1.16.2 & \texttt{UART\_Write} & Encola bytes en el buffer circular de transmisión
(128 bytes) y dispara la interrupción para iniciar el envío. \\ \hline
USW.2.1.16.4 & \texttt{UART3\_IRQHandler} & ISR: guarda los bytes recibidos y carga hasta
16 bytes en la FIFO de transmisión. \\ \hline
USW.2.1.16.5 / .6 & \texttt{UART\_CargarFifoTx}, \texttt{UART\_GuardarRecibidos} &
(Internas) Movimiento de datos entre FIFO y buffers, llamadas desde la ISR. \\ \hline
\multicolumn{3}{|l|}{\textbf{ESW.2.1.17 -- Timer}} \\ \hline
USW.2.1.17.1 & \texttt{Timer\_Init} & TIMER0 a 1\,ms por cuenta con MR0 cada 500 cuentas
e interrupción; queda detenido. \\ \hline
USW.2.1.17.14 & \texttt{Timer\_Start} & Toma la referencia del contador del sensor y
arranca TIMER0. \\ \hline
USW.2.1.17.2 & \texttt{Timer\_Read} & Devuelve el tiempo transcurrido en ms. \\ \hline
USW.2.1.17.15 & \texttt{Timer\_ConsumirVentana} & Entrega una sola vez los pulsos de la
última ventana cerrada. \\ \hline
USW.2.1.17.16 & \texttt{TIMER0\_IRQHandler} & ISR: cierra la ventana, calcula los pulsos
(diferencia módulo $2^{32}$) y reprograma MR0. \\ \hline
\multicolumn{3}{|l|}{\textbf{Soporte -- Reloj y Adquisicion\_ADC\_DMA}} \\ \hline
-- & \texttt{Reloj\_Init} & Flash a cinco ciclos, \texttt{SystemInit} (PLL0 a 100\,MHz) y
\texttt{SystemCoreClockUpdate}. \\ \hline
-- & \texttt{Adquisicion\_ADC\_DMA\_\-Init} & ADC en ráfaga sobre AD0.0 y AD0.1 y GPDMA
con lista enlazada circular hacia las muestras de cada dispositivo. \\ \hline
-- & \texttt{Adquisicion\_ADC\_DMA\_\-HayDatos} & Indica si ya se transfirió el primer par
de muestras. \\ \hline
-- & \texttt{DMA\_IRQHandler} & ISR: cuenta los pares completos. \\ \hline
\end{longtable}

\textbf{Adquisición ADC/DMA.} El ADC convierte en ráfaga los dos canales y cada fin de
conversión genera un pedido de DMA. El canal 0 del GPDMA recorre una lista enlazada
circular de dos descriptores: el primero copia AD0DR0 (potenciómetro) y el segundo AD0DR1
(LM35), cada uno a una palabra de 32 bits que pertenece al dispositivo correspondiente.
Leer siempre el registro de datos de cada canal, y nunca el registro global, fija la
asociación canal--destino aunque se pierda un pedido. Como la biblioteca no permite
cargar el primer descriptor completo, sólo los registros \texttt{DMACC0SrcAddr} y
\texttt{DMACC0Control} se escriben directamente; es el único acceso a registros fuera de
la biblioteca, junto con el de \texttt{FLASHCFG} en \texttt{Reloj\_Init}.
